\uuid{EZfx}
\chapitre{Statistique}
\niveau{L2}
\module{Analyse de données}
\sousChapitre{Tests d'hypothèses, intervalle de confiance}
\titre{Lois pour les statistiques}
\theme{loi normale, loi du chi2, loi de Student}

\auteur{}
\datecreate{2022-08-25}
\organisation{AMSCC}
\difficulte{}
\contenu{
\texte{$\mathcal N(\mu,\sigma)$ désigne la loi normale de moyenne $\mu$ et d’écart-type $\sigma$. Les quantiles sont à gauche : $q_{\nu;\gamma}$ pour la loi $\chi^2(\nu)$ ; $t_{\nu;\gamma}$ pour la loi $St(\nu)$, avec une probabilité $\gamma$ à gauche du quantile.}
\texte{Les poids $X_1,\ldots,X_{16}$ de 16 bestioles constituent un échantillon indépendant de loi $\mathcal N(100,5)$ (en grammes). On pose $\overline X=\frac1{16}\sum X_i$.}
\begin{enumerate}
\item \question{Déterminer la loi de $\overline X$.}
\reponse{Les $X_i$ sont normales et indépendantes ; leur moyenne est donc normale. Par linéarité de l'espérance,
\[
\mathbb E(\overline X)=\frac1{16}\sum_{i=1}^{16}\mathbb E(X_i)
=\frac{16\times100}{16}=100.
\]
Par indépendance,
\[
\operatorname{Var}(\overline X)=\frac1{16^2}\sum_{i=1}^{16}\operatorname{Var}(X_i)
=\frac{16\times25}{16^2}=\frac{25}{16}.
\]
Son écart-type vaut donc $\frac54$, d'où
\[
\boxed{\overline X\sim\mathcal N\left(100,\frac54\right).}
\]}
\item \question{Déterminer la loi de $Q=\frac{\sum_{i=1}^{16}(X_i-100)^2}{25}$, puis $q$ tel que $\mathbb P(Q>q)=0{,}05$.}
\reponse{On réécrit
\[
Q=\sum_{i=1}^{16}\left(\frac{X_i-100}{5}\right)^2.
\]
Les variables $\frac{X_i-100}{5}$ sont normales centrées réduites et indépendantes. Par définition de la loi du khi-deux,
\[
\boxed{Q\sim\chi^2(16).}
\]
La condition $\mathbb P(Q>q)=0{,}05$ équivaut à $\mathbb P(Q\leq q)=0{,}95$. Le seuil recherché est donc le quantile gauche d'ordre $0{,}95$ :
\[
\boxed{q=q_{16;0{,}95}\simeq26{,}296.}
\]}
\item \question{Déterminer la loi de $V=\frac{\sum_{i=1}^{16}(X_i-\overline X)^2}{25}$, puis $v$ tel que $\mathbb P(V>v)=0{,}05$.}
\reponse{Le théorème de Fisher, appliqué à cet échantillon normal indépendant, donne
\[
\boxed{V=\frac1{25}\sum_{i=1}^{16}(X_i-\overline X)^2\sim\chi^2(15).}
\]
Il donne également l'indépendance de $V$ et de $\overline X$. Le centrage sur la moyenne empirique entraîne la perte d'un degré de liberté.
Comme $\mathbb P(V>v)=0{,}05$, on cherche le quantile gauche d'ordre $0{,}95$ :
\[
\boxed{v=q_{15;0{,}95}\simeq24{,}996.}
\]}
\item \question{Déterminer la loi de $W=\frac{4\sqrt{15}(\overline X-100)}{\sqrt{\sum_{i=1}^{16}(X_i-\overline X)^2}}$, puis $w$ tel que $\mathbb P(W>w)=0{,}05$.}
\indication{Écrire $W=\frac{Z}{\sqrt{\frac{V}{15}}}$ et vérifier l'indépendance.}
\reponse{D'après la première question,
\[
Z=\frac{\overline X-100}{\frac54}=\frac{4(\overline X-100)}5\sim\mathcal N(0,1).
\]
D'autre part,
\[
\sqrt{\frac V{15}}=\frac{\sqrt{\sum_{i=1}^{16}(X_i-\overline X)^2}}{5\sqrt{15}}.
\]
Ainsi,
\[
\frac Z{\sqrt{\frac V{15}}}
=\frac{4(\overline X-100)}5\frac{5\sqrt{15}}{\sqrt{\sum_{i=1}^{16}(X_i-\overline X)^2}}=W.
\]
Le théorème de Fisher assure l'indépendance de $\overline X$ et de $V$, donc celle de $Z$ et de $V$. Les conditions de la construction de Student sont réunies :
\[
\boxed{W\sim St(15).}
\]
Enfin, $\mathbb P(W>w)=0{,}05$ donne
\[
\boxed{w=t_{15;0{,}95}\simeq1{,}75305.}
\]}
\end{enumerate}
}
